\documentclass{conjura-solution}

% amssymb is not loaded: it clashes with newtxmath (already loaded by the class) over \Bbbk;
% newtxmath already supplies \mathbb and the other symbols this document needs.
\usepackage{mathtools}

% ---------- notation ----------
\newcommand{\G}{\mathbb{G}}
\newcommand{\Gone}{\mathbb{G}_1}
\newcommand{\Gtwo}{\mathbb{G}_2}
\newcommand{\GT}{\mathbb{G}_T}
\newcommand{\Zp}{\mathbb{Z}_p}
\newcommand{\Zpx}{\mathbb{Z}_p^{\times}}
\newcommand{\vx}{\vec{x}}
\newcommand{\vy}{\vec{y}}
\newcommand{\vX}{\vec{X}}
\newcommand{\vb}{\vec{b}}
\newcommand{\Inst}{\mathsf{Inst}}
\newcommand{\Succ}{\mathsf{Succ}}
\newcommand{\Ex}{\mathbb{E}}
\newcommand{\LL}{\mathcal{L}}
\newcommand{\Span}[1]{\langle #1 \rangle}
\newcommand{\rsample}{\leftarrow_{\$}}
\newcommand{\GL}{\mathrm{GL}}
\newcommand{\AGL}{\mathrm{AGL}}
\newcommand{\bg}{\mathcal{BG}}
\newcommand{\ch}{\mathrm{ch}}
\newcommand{\Lie}{\operatorname{Lie}}
\newcommand{\Stab}{\mathrm{Stab}}
\newcommand{\agl}{\mathfrak{agl}}
\newcommand{\gl}{\mathfrak{gl}}
\newcommand{\tge}{\mathfrak{t}}
\newcommand{\vvec}{\vec{v}}   % \vv itself is already taken by newtxmath
\newcommand{\PP}{\mathbb{P}}
\newcommand{\rk}{\operatorname{rank}}
% definition/remark/proposition/theorem/corollary/lemma all come from the class

% procedure listing: framed like the class's definition boxes
\newtcolorbox{procedurebox}{enhanced jigsaw,breakable,sharp corners,
  boxrule=2pt,colframe=cjBorder,colback=white,
  left=6pt,right=6pt,top=5pt,bottom=5pt,
  before skip=9pt,after skip=9pt}

\runninghead{UBER ASSUMPTION AND RANDOM SELF-REDUCIBILITY}

\begin{document}

\cjkicker{CONJURA \textperiodcentered{} NOTE}
\cjtitle{The Uber Assumption and its Random Self-Reducibility}
\cjsubtitle{in Non-Bilinear and Type 1, 2, 3 Bilinear Groups}

\vspace{0.6em}

\begin{abstract}
We fix notation for prime-order groups without a pairing and for bilinear groups in each
of the three Galbraith--Paterson--Smart types, formally define the Uber assumption in all
four settings, and treat random self-reducibility in each. A single generic-computability
lemma isolates, per setting, the space of exponents a reduction can realize; a single
affine-reparametrization theorem then yields random self-reducibility whenever the source
data are stable under the reparametrization and the challenge polynomial transforms
covariantly. The four settings differ only through the product structure the pairing (and,
in Type 2, the homomorphism $\psi$) makes available, and this difference is exactly what
enlarges or shrinks the self-reducible subclass. We give a worked self-reducible assumption in
each setting, show that no monomial target can ever separate two settings, and exhibit a
non-monomial target that is self-reducible in Type 2 but for which, in Type 3, the affine
reduction fails at every worst-case point --- the largest the failure locus can be, since the
admissible reparametrizations form a group and the failure fraction is therefore either
$O(1/p)$ or exactly one. Because the admissibility conditions are polynomial conditions on the
reparametrization, the search for one is mechanical, and we give a decision procedure that
returns the maximal admissible group and either certifies the reduction or exhibits the locus
of worst-case points at which it fails. The procedure needs nothing bespoke: the admissible
group is the stabilizer of a subspace arrangement and of a point of a projective space, and the
failure locus is the determinantal variety of the infinitesimal action, so the whole of it is
standard computational algebra --- a Gr\"obner basis, a Jacobian kernel, an ideal of maximal
minors, a primary decomposition. A further section places the theorem against classical
random
self-reducibility. The affine condition is the one-query, linear case of the non-adaptive
condition of Feigenbaum and Fortnow; Lipton's multi-query line interpolation, which the group
encoding admits unchanged, solves the computational problem for every target here, including
the separating one, but only against a high-advantage solver and never the decisional problem.
The classical decoding machinery that would soften its threshold is not linear in the received
word and so does not survive the encoding, which is why the separation stands in the
decisional, negligible-advantage regime that cryptography works in.
\end{abstract}

\section{Introduction}

The discrete-logarithm problem is random self-reducible by the shift $g^{x}\mapsto
g^{x}g^{r}$~\cite{TW87}; the answer for a worst-case instance is recovered from an
average-case solver by subtracting $r$. This note carries the underlying idea --- injecting a
uniform reparametrization of the secret and undoing it on the answer --- across the standard
group settings, and states it uniformly. Throughout, $p$ is a prime and all groups have
order $p$. We use implicit (bracket) notation in the style of~\cite{EHKRV13}.
Section~\ref{sec:classical} relates the resulting condition to the classical form of the
same idea, in which the undoing is done by polynomial interpolation across several
correlated queries rather than in closed form on one~\cite{Lipton91,FF93}.

\section{Group settings}

\subsection{Non-bilinear groups}

\begin{definition}[prime-order group]
A \emph{(non-bilinear) group generator} outputs $\mathcal G=(p,\G,g)$ where $\G$ is a
cyclic group of prime order $p$, $g$ a generator, and the group law and equality are
computable in probabilistic polynomial time. For $a\in\Zp$ write $[a]\coloneqq g^{a}$, so
$[a]\cdot[b]=[a+b]$ and $[a]^{c}=[ca]$. There is no pairing.
\end{definition}

\subsection{Bilinear groups and the GPS taxonomy}

\begin{definition}[bilinear group]
A \emph{bilinear group generator} outputs
$\bg=(p,\Gone,\Gtwo,\GT,e,g_1,g_2)$ where $\Gone,\Gtwo,\GT$ are cyclic of prime order $p$,
$g_1,g_2$ generate $\Gone,\Gtwo$, all group laws and equalities are efficient, and
$e\colon\Gone\times\Gtwo\to\GT$ is an efficiently computable, non-degenerate bilinear map.
For $a\in\Zp$ write $[a]_1\coloneqq g_1^{a}$, $[a]_2\coloneqq g_2^{a}$, and
$[a]_T\coloneqq e(g_1,g_2)^{a}$, so that
\[
  e\bigl([a]_1,[b]_2\bigr)=[ab]_T,\qquad [a]_i\cdot[b]_i=[a+b]_i,\qquad [a]_i^{c}=[ca]_i.
\]
\end{definition}

Following Galbraith, Paterson, and Smart~\cite{GPS08}, bilinear groups come in three types,
distinguished by the efficiently computable isomorphisms between $\Gone$ and $\Gtwo$:
\begin{enumerate}[label=\textup{Type~\arabic*.},leftmargin=4.4em,itemsep=3pt]
  \item $\Gone=\Gtwo=\G$ and $e$ is symmetric, $e(u,v)=e(v,u)$. One source group.
  \item $\Gone\neq\Gtwo$, with an efficiently computable isomorphism
        $\psi\colon\Gtwo\to\Gone$ (so $\psi([a]_2)=[a]_1$), but none known in the reverse
        direction.
  \item $\Gone\neq\Gtwo$, with no efficiently computable isomorphism in either direction.
\end{enumerate}
Type 3 is the most efficient and most used in practice; Type 1 is the most permissive
algebraically. The only structural difference for what follows is which group elements can
be moved between $\Gone$ and $\Gtwo$, and hence which products the pairing can form.

\section{The generic-computability calculus}

Fix a setting and an instance consisting of encodings of finite polynomial tuples
$R\subseteq\Zp[\vX]$ in $\Gone$, $S\subseteq\Zp[\vX]$ in $\Gtwo$, and $T\subseteq\Zp[\vX]$
in $\GT$ (in a single-group setting only $R$ and, if applicable, $T$ occur). Each tuple is
assumed to contain the constant $1$, since the generators are public. For $\vX=(X_1,\dots,X_n)$
and a secret $\vx=(x_1,\dots,x_n)\in\Zp^{n}$, the instance is the collection of $[r_i(\vx)]_1$,
$[s_j(\vx)]_2$, $[t_k(\vx)]_T$.

For a group $\gamma\in\{\Gone,\Gtwo,\GT\}$ (or $\G$) let $\LL_\gamma\subseteq\Zp[\vX]$
denote the polynomials $q$ for which $[q(\vx)]_\gamma$ is computable in probabilistic
polynomial time from the instance and public parameters, uniformly in $\vx$. Write
$A\cdot B=\{ab:a\in A,\,b\in B\}$ and let $\Span{\cdot}$ be $\Zp$-linear span.

\begin{lemma}[computable exponents]
\label{lem:comp}
With the conventions above,
\[
\begin{aligned}
  \text{non-bilinear:}\quad & \LL_{\G}=\Span{R};\\[2pt]
  \text{Type 1 } (\Gone=\Gtwo=\G,\ \text{source tuple }R):\quad
      & \LL_{\G}=\Span{R},\quad \LL_{\GT}=\Span{R\cdot R\ \cup\ T};\\[2pt]
  \text{Type 2 } (\psi\colon\Gtwo\to\Gone):\quad
      & \LL_{1}=\Span{R\cup S},\quad \LL_{2}=\Span{S},\\
      &\quad \LL_{T}=\Span{(R\cup S)\cdot S\ \cup\ T};\\[2pt]
  \text{Type 3:}\quad
      & \LL_{1}=\Span{R},\quad \LL_{2}=\Span{S},\quad
        \LL_{T}=\Span{R\cdot S\ \cup\ T}.
\end{aligned}
\]
The inclusions ``$\supseteq$'' hold unconditionally: every listed polynomial's encoding is
computable by group operations and pairings. The reverse inclusions hold in the generic
bilinear group model, where these are exactly the computable encodings~\cite{BBG05,Boyen08}.
\end{lemma}

\begin{proof}[Justification]
The group law realizes all $\Zp$-linear combinations, giving $\Span{\cdot}$ in each group.
The pairing sends a computable $\Gone$-exponent $a\in\LL_1$ and a computable
$\Gtwo$-exponent $b\in\LL_2$ to $ab\in\LL_T$; pairing against the generator ($b=1$ or
$a=1$) shows $\LL_1,\LL_2\subseteq\LL_T$. In Type 1 both pairing arguments range over
$\Span{R}$, producing all products $R\cdot R$. In Type 2, $\psi$ moves each $\Gtwo$
encoding into $\Gone$, so $\LL_1$ gains $S$, and the pairable products become
$\LL_1\cdot\LL_2=(R\cup S)\cdot S$; the missing products $R\cdot R$ would require an
element of $\Gtwo$ with an $R$-exponent, which $\psi$ cannot supply. In Type 3 neither
group maps to the other, so $\LL_1=\Span{R}$, $\LL_2=\Span{S}$, and only cross products
$R\cdot S$ arise. The generic-model converse is the master lemma
of~\cite{BBG05,Boyen08}.
\end{proof}

The single distinction across the four rows is the set of $\GT$-products: none (no pairing),
$R\cdot R$ (Type 1), $R\cdot S\cup S\cdot S$ (Type 2), $R\cdot S$ (Type 3). Everything below
is a consequence of this.

\section{The Uber assumption}

\begin{definition}[bilinear Uber problem]
\label{def:uber}
Fix a bilinear setting, source tuples $R$ (in $\Gone$), $S$ (in $\Gtwo$), $T$ (in $\GT$),
each containing $1$, a challenge polynomial $f\in\Zp[\vX]$, and a challenge group
$\ch\in\{\Gone,\Gtwo,\GT\}$; write $\LL_\star\coloneqq\LL_{\ch}$ (Lemma~\ref{lem:comp}).
For a secret $\vx\in\Zp^{n}$ the instance is
\[
  \Inst(\vx)=\bigl(\{[r_i(\vx)]_1\}_i,\ \{[s_j(\vx)]_2\}_j,\ \{[t_k(\vx)]_T\}_k\bigr).
\]
\begin{itemize}[leftmargin=1.4em,itemsep=2pt]
  \item \emph{Computational} $(R,S,T,f;\ch)$-Uber: given $\Inst(\vx)$, output
        $[f(\vx)]_\star$.
  \item \emph{Decisional} $(R,S,T,f;\ch)$-Uber: given $\Inst(\vx)$ and $Z$, decide whether
        $Z=[f(\vx)]_\star$ (real) or $Z\rsample\ch$ (random).
\end{itemize}
The \emph{average case} draws $\vx\rsample\Zp^{n}$; the \emph{worst case} quantifies over
all $\vx$. For a decision algorithm $\mathcal A$ put
\[
  \delta(\mathcal A;\vx)=\Pr\bigl[\mathcal A(\Inst(\vx),[f(\vx)]_\star)=1\bigr]
  -\Pr\bigl[\mathcal A(\Inst(\vx),U)=1\bigr],\quad U\rsample\ch,
\]
and $\delta_{\mathrm{avg}}(\mathcal A)=\Ex_{\vx\rsample\Zp^{n}}[\delta(\mathcal A;\vx)]$;
for a search algorithm put $\Succ(\mathcal A;\vx)=\Pr[\mathcal A(\Inst(\vx))=[f(\vx)]_\star]$
and $\Succ_{\mathrm{avg}}(\mathcal A)=\Ex_{\vx}[\Succ(\mathcal A;\vx)]$. The assumption
asserts that $\delta_{\mathrm{avg}}$ (resp.\ $\Succ_{\mathrm{avg}}$) is negligible for all
probabilistic polynomial-time $\mathcal A$.
\end{definition}

The four settings are Definition~\ref{def:uber} with the computable spaces of
Lemma~\ref{lem:comp}. We record each explicitly.

\begin{definition}[non-bilinear Uber]
Setting $\mathcal G=(p,\G,g)$; source tuple $R$ (in $\G$, containing $1$); challenge
$f\in\Zp[\vX]$ in $\G$. Instance $\Inst(\vx)=\{[r_i(\vx)]\}_i$; goal, compute or distinguish
$[f(\vx)]$. Computable space $\LL_\star=\LL_{\G}=\Span{R}$; no pairing, hence no products.
\end{definition}

\begin{definition}[Type-1 Uber]
Setting $\bg$ of Type 1, $\Gone=\Gtwo=\G$; source tuple $R$ (in $\G$), target tuple $T$
(in $\GT$), each containing $1$; challenge $f$ in $\ch\in\{\G,\GT\}$. Computable spaces
$\LL_{\G}=\Span{R}$ and $\LL_{\GT}=\Span{R\cdot R\cup T}$.
\end{definition}

\begin{definition}[Type-2 Uber]
Setting $\bg$ of Type 2 with $\psi\colon\Gtwo\to\Gone$; source tuples $R$ (in $\Gone$),
$S$ (in $\Gtwo$), $T$ (in $\GT$); challenge $f$ in $\ch\in\{\Gone,\Gtwo,\GT\}$. Computable
spaces $\LL_1=\Span{R\cup S}$, $\LL_2=\Span{S}$,
$\LL_T=\Span{(R\cup S)\cdot S\cup T}$.
\end{definition}

\begin{definition}[Type-3 Uber]
Setting $\bg$ of Type 3; source tuples $R$ (in $\Gone$), $S$ (in $\Gtwo$), $T$ (in $\GT$);
challenge $f$ in $\ch\in\{\Gone,\Gtwo,\GT\}$. Computable spaces $\LL_1=\Span{R}$,
$\LL_2=\Span{S}$, $\LL_T=\Span{R\cdot S\cup T}$.
\end{definition}

\begin{remark}[hardness]
The master theorem of~\cite{BBG05,Boyen08} states, informally, that the decisional
$(R,S,T,f;\ch)$-Uber problem is hard in the generic bilinear group model if and only if
$f\notin\LL_\star$, that is, the challenge exponent is linearly independent of the
computable ones. Random self-reducibility is a separate property and does not require
hardness; the two interact through Lemma~\ref{lem:comp}, as discussed in
Section~\ref{sec:discussion}.
\end{remark}

\section{An affine random-self-reducibility theorem}

Let $\AGL_n(\Zp)$ be the affine group of maps $\tau_{A,\vb}\colon\vx\mapsto A\vx+\vb$ with
$A\in\GL_n(\Zp)$, acting on polynomials by $(q\circ\tau)(\vX)=q(A\vX+\vb)$. The following
theorem is setting-agnostic: it refers only to the computable spaces of
Lemma~\ref{lem:comp}.

\begin{theorem}[master reduction]
\label{thm:master}
Consider any of the four settings, an instance with source tuples $R,S,T$ and challenge
$(f;\ch)$, and computable spaces $\{\LL_\gamma\}$ with $\LL_\star=\LL_{\ch}$. Suppose there
is an efficiently sampleable distribution $\mathcal D$ on a subgroup
$\mathcal T\le\AGL_n(\Zp)$ such that, for every $\tau\in\mathcal T$,
\begin{enumerate}[label=\textup{(C\arabic*)},leftmargin=3.4em,itemsep=3pt]
  \item $r_i\circ\tau\in\LL_1$, $s_j\circ\tau\in\LL_2$, and $t_k\circ\tau\in\LL_T$ for all
        $i,j,k$ (with $\LL_1=\LL_{\G}$ in the single-group settings);
  \item $f\circ\tau=\lambda_\tau f+h_\tau$ for a scalar $\lambda_\tau\in\Zpx$ and a
        polynomial $h_\tau\in\LL_\star$, both computable from $\tau$ and the parameters;
\end{enumerate}
and, moreover,
\begin{enumerate}[label=\textup{(C\arabic*)},leftmargin=3.4em,start=3]
  \item $\tau(\vx)$ is uniform on $\Zp^{n}$ for every fixed $\vx$ when $\tau\leftarrow\mathcal D$.
\end{enumerate}
Then the computational and decisional $(R,S,T,f;\ch)$-Uber problems are random
self-reducible: there is a probabilistic polynomial-time reduction $\mathcal B^{\mathcal A}$,
using one draw from $\mathcal D$ and one oracle call, with
$\Succ(\mathcal B;\vx_0)=\Succ_{\mathrm{avg}}(\mathcal A)$ and
$\delta(\mathcal B;\vx_0)=\delta_{\mathrm{avg}}(\mathcal A)$ for every $\vx_0\in\Zp^{n}$.
\end{theorem}

\begin{proof}
On input $\Inst(\vx_0)$ (and a challenge $Z$ in the decisional case) the reduction proceeds
as follows.

\emph{Reparametrize.} Draw $\tau\leftarrow\mathcal D$ and set $\vx'\coloneqq\tau(\vx_0)$,
which is unknown to $\mathcal B$ but uniform on $\Zp^{n}$ by (C3).

\emph{Rebuild the instance.} By (C1), each rebuilt exponent lies in the computable space of
its group, so by the ``$\supseteq$'' part of Lemma~\ref{lem:comp} its encoding is
computable from $\Inst(\vx_0)$. Concretely, writing $r_i\circ\tau=\sum_a c_a\,g_a$ as a
$\Zp$-combination of the generators of $\LL_1$, one has
$[r_i(\vx')]_1=[(r_i\circ\tau)(\vx_0)]_1=\prod_a[g_a(\vx_0)]_1^{c_a}$, and similarly for
$S$ (using $\psi$ in Type 2 to realize any $S$-exponent in $\Gone$) and for $T$ (using
pairings $e([\cdot]_1,[\cdot]_2)$ to realize products). This yields $\Inst(\vx')$,
distributed exactly as a fresh average-case instance.

\emph{Transport the challenge.} By (C2) compute $\lambda_\tau,h_\tau$ and the element
$[h_\tau(\vx_0)]_\star$ (computable since $h_\tau\in\LL_\star$).

\emph{Decisional case.} Set $Z'\coloneqq Z^{\lambda_\tau}\cdot[h_\tau(\vx_0)]_\star$ and
output $\mathcal A(\Inst(\vx'),Z')$. If $Z=[f(\vx_0)]_\star$ then
\[
  Z'=[\lambda_\tau f(\vx_0)+h_\tau(\vx_0)]_\star=[(f\circ\tau)(\vx_0)]_\star=[f(\vx')]_\star,
\]
the real challenge for $\vx'$. If $Z$ is uniform on $\ch$, then $Z^{\lambda_\tau}$ is
uniform because $z\mapsto z^{\lambda_\tau}$ is a bijection of the order-$p$ cyclic group
$\ch$ (as $\gcd(\lambda_\tau,p)=1$), so $Z'$ is uniform and independent of $\Inst(\vx')$.
Writing $p_1,p_0$ for $\mathcal A$'s acceptance probabilities on real and random challenges,
\[
  \delta(\mathcal B;\vx_0)
  =\Ex_{\tau}[p_1(\tau(\vx_0))]-\Ex_{\tau}[p_0(\tau(\vx_0))]
  \overset{\text{(C3)}}{=}\Ex_{\vx'}[p_1]-\Ex_{\vx'}[p_0]
  =\delta_{\mathrm{avg}}(\mathcal A).
\]

\emph{Search case.} Run $W\leftarrow\mathcal A(\Inst(\vx'))$. On success
$W=[f(\vx')]_\star=[\lambda_\tau f(\vx_0)+h_\tau(\vx_0)]_\star$, so output
$Y\coloneqq(W\cdot[h_\tau(\vx_0)]_\star^{-1})^{\lambda_\tau^{-1}\bmod p}=[f(\vx_0)]_\star$.
Since $\vx'$ is uniform, the success probability equals $\Succ_{\mathrm{avg}}(\mathcal A)$.
\end{proof}

\begin{remark}[relaxing (C3)]
\label{rem:c3}
When $\mathcal D$ makes $\tau(\vx)$ uniform only on a subset $\Omega\subseteq\Zp^{n}$ (for
instance $\Zpx\times\Zp$ under a multiplicative reparametrization of one coordinate), the
transfer holds with the average taken over $\Omega$. If the complement $\Zp^{n}\setminus
\Omega$ is a proper affine subvariety on which $f(\vx)$ takes a value the reduction can
compute directly from $\Inst(\vx_0)$ (a common case is $f(\vx)=0$ there), the reduction
detects that case by an equality test on the given encodings and answers it without an
oracle call, recovering full worst-case coverage.
\end{remark}

\section{Random self-reducibility: non-bilinear groups}
\label{sec:nonbilinear}

Here $\LL_\star=\LL_{\G}=\Span{R}$ and no pairing exists, so (C2) demands the covariance
defect to be a $\Zp$-combination of the source exponents themselves. The motivating
scalar case is discrete log: to solve $g^{x}$ query the solver on $g^{x+r}$ and subtract
$r$~\cite{TW87}; its group-valued cousins fit Definition~\ref{def:uber}.

\begin{proposition}[CDH and DDH]
\label{prop:cdh}
Let $R=(1,X_1,X_2)$, so $\LL_{\G}=\Span{1,X_1,X_2}$, and $f=X_1X_2$ with challenge in $\G$.
The translation group $\mathcal T=\{\vx\mapsto\vx+\vb\}$ with $\vb\rsample\Zp^{2}$ satisfies
\textup{(C1)--(C3)}, so computational CDH and decisional DDH are random self-reducible.
\end{proposition}
\begin{proof}
$X_i\circ\tau=X_i+b_i\in\LL_{\G}$ gives (C1). For (C2),
$f\circ\tau=(X_1+b_1)(X_2+b_2)=f+(b_2X_1+b_1X_2+b_1b_2)$, so $\lambda_\tau=1$ and
$h_\tau=b_2X_1+b_1X_2+b_1b_2\in\Span{1,X_1,X_2}=\LL_\star$. Translation gives exact
uniformity, (C3). Explicitly the reduction returns
$[f(\vx_0)]=W\cdot[x_1]^{-b_2}[x_2]^{-b_1}[1]^{-b_1b_2}$ in the search case.
\end{proof}

With $R=(1,X_1,\dots,X_n)$ we have $\LL_\star=V_{\le 1}$ (affine polynomials), and
translating a degree-$d$ monomial leaves monomials of degree at most $d-1$; hence (C2)
holds under translation exactly when $\deg f\le 2$. Without a pairing, degree $2$ is the
ceiling: a product target is the most that is randomizable, and it is.

\section{Random self-reducibility: Type 1}

Here $\LL_\star=\LL_{\GT}=\Span{R\cdot R\cup T}$: the pairing supplies all products of source
exponents, so the covariance defect may use any degree-$2$ combination of them.

\begin{proposition}[BDDH]
\label{prop:bddh}
Let $R=(1,X_1,X_2,X_3)$, $T=(1)$, so $\LL_{\GT}=V_{\le 2}$, and $f=X_1X_2X_3$ with challenge
in $\GT$. Translation satisfies \textup{(C1)--(C3)}, so decisional and computational BDDH are
random self-reducible.
\end{proposition}
\begin{proof}
$X_i\circ\tau=X_i+b_i\in\LL_{\G}=\Span{R}$ gives (C1). Expanding,
\[
  f\circ\tau=X_1X_2X_3+\bigl[b_3X_1X_2+b_2X_1X_3+b_1X_2X_3+(\deg\le 1)\bigr],
\]
so $\lambda_\tau=1$ and $h_\tau$ has total degree at most $2$, hence lies in
$V_{\le 2}=\LL_\star$; each of its monomials is realized by a single pairing, for instance
$[x_1x_2]_T=e([x_1]_1,[x_2]_2)$. Translation gives (C3).
\end{proof}

With coordinate source $R=(1,X_1,\dots,X_n)$ we get $\LL_\star=V_{\le 2}$, so the same
degree count gives translation-randomizability for every $f$ with $\deg f\le 3$. The
pairing has raised the ceiling from $2$ (no pairing) to $3$.

\section{Random self-reducibility: Type 2}

Here $\LL_1=\Span{R\cup S}$ and $\LL_T=\Span{(R\cup S)\cdot S\cup T}$: relative to Type 3
the homomorphism $\psi$ adds the right-hand products $S\cdot S$ to the target space. This
extra room is what makes the following target randomizable.

\begin{proposition}
\label{prop:type2}
In a Type-2 group let $R=(1,X_1)$ in $\Gone$, $S=(1,X_2)$ in $\Gtwo$, $T=(1)$, and
$f=X_1X_2^{2}$ with challenge in $\GT$. Then $f\notin\LL_\star$ (so the assumption is
generically hard), and translation satisfies \textup{(C1)--(C3)}; the problem is random
self-reducible.
\end{proposition}
\begin{proof}
Here $\LL_T=\Span{(R\cup S)\cdot S}=\Span{1,X_1,X_2,X_1X_2,X_2^{2}}$, and $X_1X_2^{2}$ is
not in this span, giving generic hardness. Under $\tau\colon\vx\mapsto\vx+\vb$ one has
$X_1\circ\tau=X_1+b_1\in\LL_1=\Span{1,X_1}$ and $X_2\circ\tau=X_2+b_2\in\LL_2=\Span{1,X_2}$,
which is (C1). Expanding,
\[
  f\circ\tau=(X_1+b_1)(X_2+b_2)^{2}
  =f+\underbrace{2b_2X_1X_2+b_2^{2}X_1+b_1X_2^{2}+2b_1b_2X_2+b_1b_2^{2}}_{h_\tau},
\]
so $\lambda_\tau=1$ and every monomial of $h_\tau$ lies in $\LL_\star$. The term $X_2^{2}$
is realized using $\psi$: $[x_2^{2}]_T=e\bigl(\psi([x_2]_2),[x_2]_2\bigr)$. Translation gives
(C3).
\end{proof}

The point is precisely the $X_2^{2}$ term of $h_\tau$: it is a product of two $\Gtwo$
exponents, which the pairing can form only after $\psi$ has moved one of them into $\Gone$.
In Type 1 this term is free; in Type 3 it is unavailable, as the next section shows.

\section{Random self-reducibility: Type 3}
\label{sec:type3}

Here $\LL_1=\Span{R}$, $\LL_2=\Span{S}$, and $\LL_T=\Span{R\cdot S\cup T}$: the target space
contains only cross products. This is the most restrictive setting, and reparametrizations
must respect which coordinates are reachable in the challenge group.

\begin{proposition}[co-CDH]
\label{prop:cocdh}
In a Type-3 group let $R=(1,X_1)$ in $\Gone$, $S=(1,X_2)$ in $\Gtwo$, and $f=X_1X_2$ with
challenge in $\Gone$. Then $f\notin\LL_\star$ (generically hard), and the reparametrization
$\tau\colon(x_1,x_2)\mapsto(\alpha x_1,\ x_2+b_2)$ with $\alpha\rsample\Zpx$,
$b_2\rsample\Zp$ satisfies \textup{(C1),(C2)} and makes $\tau(\vx)$ uniform on
$\Zpx\times\Zp$; the problem is random self-reducible, with the slice $x_1=0$ handled
directly.
\end{proposition}
\begin{proof}
Here $\LL_\star=\LL_1=\Span{1,X_1}$, which contains no $\Gtwo$ variable; and $X_1X_2\notin
\Span{1,X_1}$, giving hardness. For (C1), $X_1\circ\tau=\alpha X_1\in\LL_1$ and
$X_2\circ\tau=X_2+b_2\in\LL_2=\Span{1,X_2}$. For (C2),
$f\circ\tau=\alpha X_1(X_2+b_2)=\alpha f+\alpha b_2X_1$, so $\lambda_\tau=\alpha\in\Zpx$ and
$h_\tau=\alpha b_2X_1\in\LL_\star$. The image $\tau(\vx)=(\alpha x_1,x_2+b_2)$ is uniform on
$\Zpx\times\Zp$ for every fixed $x_1\neq 0$. If $x_1=0$ then $f(\vx)=0$ and the answer is
the identity of $\Gone$; the reduction detects this by testing $[x_1]_1=[0]_1$ and outputs
the identity, per Remark~\ref{rem:c3}. A multiplicative reparametrization of $x_1$ is forced
because an additive shift $x_1\mapsto x_1+b_1$ would put $b_1X_2$ into $h_\tau$, and
$X_2\notin\LL_1$ in Type 3.
\end{proof}

\subsection{Monomial targets never separate}

It is tempting to look for a separation at the target $f=X_1X_2^{2}$ of
Proposition~\ref{prop:type2}, whose defect under translation contains the $X_2^{2}$ that
Type 3 cannot form. That attempt fails, and it fails for a reason that rules out an entire
class of candidates at once.

\begin{proposition}[monomial targets are always self-reducible]
\label{prop:nomono}
Let $R,S,T$ consist of monomials, each containing $1$, and let $f=X_1^{i}X_2^{j}$ with
$i,j\ge 1$ be a monomial with $f\notin\LL_\star$. Then in every one of the four settings the
diagonal torus $\mathcal T=\{\tau_{a,d}\colon\vx\mapsto(ax_1,dx_2):a,d\in\Zpx\}$, with $a$ and
$d$ drawn uniformly from $\Zpx$, satisfies
\textup{(C1)} and \textup{(C2)}, and Theorem~\ref{thm:master} applies in the relaxed form of
Remark~\ref{rem:c3}. No such $f$ separates two settings.
\end{proposition}
\begin{proof}
A monomial composed with a coordinate scaling is a scalar multiple of itself, so
$m\circ\tau_{a,d}=a^{k}d^{l}m$ for $m=X_1^{k}X_2^{l}$. Hence every source monomial stays in
the span of itself, giving (C1) in any setting, and $f\circ\tau_{a,d}=a^{i}d^{j}f$, giving
(C2) with $\lambda_\tau=a^{i}d^{j}\in\Zpx$ and $h_\tau=0$ --- the defect vanishes identically,
so no uncomputable monomial can appear in it. The image $\tau_{a,d}(\vx)$ is uniform on
$\Zpx\times\Zpx$ whenever $x_1,x_2\neq 0$, and on the complement
$\{x_1=0\}\cup\{x_2=0\}$ one has $f(\vx)=0$ because $i,j\ge 1$, so the answer is the identity
and the reduction detects the case by testing $[x_1]_1$ and $[x_2]_2$ against the identity.
Remark~\ref{rem:c3} then gives full worst-case coverage.
\end{proof}

In particular $f=X_1X_2^{2}$ is random self-reducible in Type 3 as well as in Type 2, by
$\tau\colon\vx\mapsto(\alpha x_1,\,x_2+b)$: here $f\circ\tau=\alpha f+2\alpha b\,X_1X_2+
\alpha b^{2}X_1$, whose defect monomials $X_1X_2$ and $X_1$ are both cross products $R\cdot S$
and hence in $\LL_T$ even in Type 3. This is the same shape of reparametrization that
Proposition~\ref{prop:cocdh} forces for co-CDH. A separating target must therefore be
non-monomial, and must defeat the torus.

\subsection{A separating target}

\begin{proposition}[a Type-2 versus Type-3 separation]
\label{prop:sep}
Let $p\ge 5$, $R=(1,X_1)$ in $\Gone$, $S=(1,X_2)$ in $\Gtwo$, $T=(1)$, and
\[
  f=X_1^{2}+X_1X_2^{2}+X_2^{3}
\]
with challenge in $\GT$. Then $f\notin\LL_T$ in Type 2 and in Type 3, so the assumption is
generically hard in both. In Type 2 the admissible reparametrizations act transitively and
uniformly on $\Zp^{2}$, so Theorem~\ref{thm:master} applies with exact advantage transfer. In
Type 3 the admissible group is
\[
  \mathcal T_3=\{\tau_c\colon\vx\mapsto(x_1-3c,\ x_2+c)\ :\ c\in\Zp\}\cong(\Zp,+),
\]
of order $p$; its orbits are the $p$ lines $x_1+3x_2=k$, each of size $p$; and on no such line
is $f$ computable from the instance. Hence \textup{(C3)} fails at \emph{every}
$\vx_0\in\Zp^{2}$, even in the relaxed form of Remark~\ref{rem:c3}.
\end{proposition}
\begin{proof}
\emph{Hardness.} In Type 3, $\LL_T=\Span{R\cdot S}=\Span{1,X_1,X_2,X_1X_2}$; in Type 2,
$\LL_T=\Span{(R\cup S)\cdot S}=\Span{1,X_1,X_2,X_1X_2,X_2^{2}}$. Neither contains $X_1^{2}$,
so $f\notin\LL_T$ in either.

\emph{Type 3.} By (C1), $\tau_1=aX_1+c_1$ and $\tau_2=dX_2+c_2$ with $a,d\in\Zpx$. The
monomials of $f\circ\tau$ lying outside $\LL_T$ are $X_1^{2}$, $X_1X_2^{2}$, $X_2^{3}$ and
$X_2^{2}$; the first three occur in $f$ and the fourth does not. Reading off coefficients,
\[
  [X_1^{2}]:\ a^{2}=\lambda_\tau,\qquad
  [X_1X_2^{2}]:\ ad^{2}=\lambda_\tau,\qquad
  [X_2^{3}]:\ d^{3}=\lambda_\tau .
\]
From the first two, $a=d^{2}$; from the last two, $a=d$; hence $d^{2}=d$ and, as
$d\in\Zpx$, $d=1$, so $a=1$ and $\lambda_\tau=1$. The pure $X_2^{2}$ coefficient is
$c_1+3c_2$ and must vanish, giving $c_1=-3c_2$. Writing $c\coloneqq c_2$ yields exactly
$\mathcal T_3$, and one checks $h_{\tau_c}=2c\,X_1X_2+(c^{2}-6c)X_1-3c^{2}X_2+(9c^{2}-2c^{3})
\in\LL_T$, so (C2) does hold on $\mathcal T_3$. Since $x_1+3x_2$ is invariant under $\tau_c$,
the orbit of $\vx_0$ is the line $\ell_k=\{x_1+3x_2=k\}$ through it, of size $p$ out of
$p^{2}$.

\emph{No escape on an orbit.} Parametrize $\ell_k$ by $s\coloneqq x_2$, so $x_1=k-3s$. There
\[
  f|_{\ell_k}=(k-3s)^{2}+(k-3s)s^{2}+s^{3}=-2s^{3}+(k+9)s^{2}-6ks+k^{2},
\]
while $\LL_T$ restricts to $\Span{1,\,k-3s,\,s,\,(k-3s)s}=\Span{1,s,s^{2}}$. As $p\ge 5$ the
coefficient $-2$ is nonzero and $s^{3}$ is not a $\Zp$-combination of $1,s,s^{2}$ as a
function on $\Zp$, so $f|_{\ell_k}\notin\LL_T|_{\ell_k}$ for every $k$. The escape clause of
Remark~\ref{rem:c3} is therefore unavailable on any orbit.

\emph{Type 2.} Now $\LL_1=\Span{1,X_1,X_2}$, so (C1) permits $\tau_1=a_{11}X_1+a_{12}X_2+c_1$,
and $X_2^{2}\in\LL_T$ removes the $X_2^{2}$ constraint. The surviving conditions are
$a_{11}^{2}=\lambda_\tau$, $a_{11}d^{2}=\lambda_\tau$ and $a_{12}d^{2}+d^{3}=\lambda_\tau$,
giving $a_{11}=d^{2}$ and $a_{12}=d^{2}-d$, with $c_1,c_2$ free. Then
$\tau(\vx)=(d^{2}x_1+(d^{2}-d)x_2+c_1,\ dx_2+c_2)$, and as $c_1,c_2$ range over $\Zp$ the two
coordinates range independently over all of $\Zp$: $\tau(\vx)$ is exactly uniform on $\Zp^{2}$
for every fixed $\vx$, which is (C3) unrelaxed.
\end{proof}

Two features of Type 2 drive the separation, and the proof uses both: $\psi$ places $X_2^{2}$
in $\LL_T$, which removes the constraint that ties $c_1$ to $c_2$ in Type 3, and it places
$X_2$ in $\LL_1$, which lets $\tau_1$ carry the shear term $a_{12}X_2$. Neither is available
in Type 3, and either alone would leave a smaller admissible group.

On the other side, all three terms of $f$ are needed, and the coefficient conditions
$a^{2}=ad^{2}=d^{3}=\lambda_\tau$ are what make them work together: any two of the three force
only a single relation between $a$ and $d$, leaving a one-parameter torus, whereas all three
force $d=1$ and kill the torus that Proposition~\ref{prop:nomono} would otherwise use to
rescue the Type-3 reduction. Dropping a term accordingly costs the separation, but in
different ways. Without $X_1^{2}$ one has $a=d$; the torus survives, and the failure locus
shrinks to the single line $x_1+3x_2=0$, an $O(1/p)$ fraction. Without $X_1X_2^{2}$ one has
$a^{2}=d^{3}$, and without $X_2^{3}$ one has $a=d^{2}$; in both cases the surviving
one-parameter torus leaves Type-3 orbits of size $\Theta(p^{2})$, so the loss is a constant
factor rather than a break --- Type 2 retains full coverage in each case, so what fails is the
Type-3 half of the separation, not the Type-2 half.

The hypothesis $p\ge 5$ is not cosmetic. The proof needs $s^{3}$ to be independent of
$1,s,s^{2}$ \emph{as a function} on $\Zp$, which holds exactly when $p>3$. At $p=3$ Fermat
gives $s^{3}=s$, so $f|_{\ell_k}=ks^{2}+s+k^{2}$ lies in $\LL_T|_{\ell_k}$ and the escape
clause of Remark~\ref{rem:c3} becomes available on every orbit; at $p=2$ the assumption is
vacuous, since $X_1^{2},X_2^{2}$ and $X_2^{3}$ collapse to $X_1,X_2,X_2$ and $f$ agrees
pointwise with $X_1+X_1X_2+X_2\in\LL_T$. Both failures are of hardness or of the escape
clause, not of the group computation, which is characteristic-free.

\begin{remark}[the failure is as large as it can be]
\label{rem:dichotomy}
The set of $\tau$ satisfying (C1) and (C2) is closed under composition and inversion, so it is
a group and the orbits of $\Zp^{n}$ under it partition the space. Orbit size is constant on
orbits, so the set of $\vx_0$ at which (C3) fails is a union of orbits. If it is not all of
$\Zp^{n}$, it omits some orbit of size at least $p^{n}-O(p)$ and therefore has size $O(p)$.
The failure fraction is thus either $O(1/p)$ or exactly $1$, with nothing in between, and
Proposition~\ref{prop:sep} realizes the second alternative.
\end{remark}

Proposition~\ref{prop:sep} rules out the single-reparametrization technique of
Theorem~\ref{thm:master}, which is the tool behind every positive result above. It does not
rule out every reduction: Section~\ref{sec:classical} shows that a multi-query interpolation
reduction still solves the \emph{computational} problem, at high solver advantage, for this
target as for any other with these source tuples. The Type-2 versus Type-3 gap established
here is accordingly a gap for decisional reductions and for reductions that must work at
negligible advantage. Ruling out every decisional reduction would additionally require a
lower-bound argument such as a meta-reduction.

\section{Relation to classical random self-reducibility}
\label{sec:classical}

Random self-reducibility long predates its use on group-based assumptions, and the classical
technique differs from Theorem~\ref{thm:master} in one respect that decides exactly what
Proposition~\ref{prop:sep} does and does not rule out. This section isolates the difference
and delimits the reach of each technique.

\subsection{The multi-query condition}

Condition (C2) asks a single reparametrized challenge to be a scalar multiple of the original
up to a computable defect. Feigenbaum and Fortnow~\cite{FF93} call a problem non-adaptively
random self-reducible when $f(\vx_0)$ is recoverable by an arbitrary polynomial-time function
of the answers on several identically distributed queries; see~\cite{BT06} for the place of
that notion in average-case complexity. The specialization in which the recovery is
$\Zp$-linear is the one this setting supports:
\begin{enumerate}[label=\textup{(C2$'$)},leftmargin=3.9em]
  \item there are $\tau_1,\dots,\tau_m$, jointly sampled so that each $\tau_i$ is marginally
        distributed as in (C3), together with coefficients $c_1,\dots,c_m\in\Zp$ and a
        polynomial $h\in\LL_\star$, all computable from the sampled data, such that
        \[
          f=\sum_{i=1}^{m}c_i\,(f\circ\tau_i)+h.
        \]
\end{enumerate}
Condition (C2) is the case $m=1$ with $c_1\neq 0$: from $f\circ\tau=\lambda_\tau f+h_\tau$ one
gets $f=\lambda_\tau^{-1}(f\circ\tau)-\lambda_\tau^{-1}h_\tau$, which is (C2$'$) with
$c_1=\lambda_\tau^{-1}$ and $h=-\lambda_\tau^{-1}h_\tau$.

Given (C1) for each $\tau_i$, condition (C2$'$) suffices for the \emph{computational} problem:
the reduction rebuilds $\Inst(\tau_i(\vx_0))$ for each $i$ as in the proof of
Theorem~\ref{thm:master}, queries the solver on all $m$ of them to get $W_1,\dots,W_m$, and
returns
\[
  \prod_{i=1}^{m}W_i^{\,c_i}\cdot[h(\vx_0)]_\star
  =\Bigl[\sum_{i=1}^{m}c_i\,f(\tau_i(\vx_0))+h(\vx_0)\Bigr]_\star=[f(\vx_0)]_\star .
\]
The step that makes this go through is that a $\Zp$-linear combination of \emph{exponents} is
realized by group operations on their encodings, so linear post-processing of the solver's
answers is available to the reduction for free. Encodings are transparent to it: the map
$\vx\mapsto[f(\vx)]_\star$ is a Reed--Muller codeword read through the encoding, and a linear
local-decoding procedure for that code never notices the encoding is there. This is why the
linear part of the classical toolkit --- Lipton's line restriction~\cite{Lipton91}, and the
degree-$k$ curves with which Beaver and Feigenbaum~\cite{BF90} make the queries $k$-wise
independent --- carries over to computational Uber problems unchanged.

\subsection{Interpolation, and what it costs}

The classical instantiation restricts $f$ to a random line through $\vx_0$ and interpolates.
It needs nothing of $f$ beyond its degree, and nothing of the setting beyond (C1) for
translations.

\begin{proposition}[line interpolation]
\label{prop:interp}
Suppose $R$, $S$ and $T$ are stable under translation, in the sense that $r\circ\tau_{\vb}\in
\LL_1$, $s\circ\tau_{\vb}\in\LL_2$ and $t\circ\tau_{\vb}\in\LL_T$ for every $\vb\in\Zp^{n}$ and
every $r\in R$, $s\in S$, $t\in T$, where $\tau_{\vb}\colon\vx\mapsto\vx+\vb$. Let
$d\coloneqq\deg f$ and suppose $d\le p-2$. Then the \emph{computational}
$(R,S,T,f;\ch)$-Uber problem is random self-reducible by a non-adaptive $(d+1)$-query
reduction $\mathcal B^{\mathcal A}$ with
\[
  \Succ(\mathcal B;\vx_0)\ \ge\ 1-(d+1)\bigl(1-\Succ_{\mathrm{avg}}(\mathcal A)\bigr)
  \qquad\text{for every }\vx_0\in\Zp^{n}.
\]
\end{proposition}
\begin{proof}
Draw $\vb\rsample\Zp^{n}$ and fix $d+1$ distinct nonzero $t_1,\dots,t_{d+1}\in\Zp$, which exist
as $d+1\le p-1$. Put $\tau_i\coloneqq\tau_{t_i\vb}$. Each $\tau_i$ is a translation, so (C1)
holds by hypothesis and $\Inst(\vx_0+t_i\vb)$ is computable from $\Inst(\vx_0)$; and since
$t_i\neq 0$ and $\vb$ is uniform, $\vx_0+t_i\vb$ is uniform on $\Zp^{n}$, so each query is
marginally a fresh average-case instance. The $d+1$ queries are not jointly independent: they
lie on one line.

Put $g(t)\coloneqq f(\vx_0+t\vb)$, a polynomial in $t$ of degree at most $d$, and let
$c_i\coloneqq\prod_{j\neq i}(-t_j)(t_i-t_j)^{-1}$ be the Lagrange coefficients at $0$, so that
$q(0)=\sum_i c_i q(t_i)$ for every $q$ of degree at most $d$. This is (C2$'$) with $m=d+1$ and
$h=0$. Query $W_i\leftarrow\mathcal A(\Inst(\vx_0+t_i\vb))$ and output
$Y\coloneqq\prod_i W_i^{\,c_i}$; if every $W_i=[g(t_i)]_\star$ then
$Y=[\sum_i c_i g(t_i)]_\star=[g(0)]_\star=[f(\vx_0)]_\star$. A union bound over the $d+1$
queries gives the stated bound.
\end{proof}

The source tuples of Propositions~\ref{prop:cocdh}, \ref{prop:nomono} and~\ref{prop:sep} are
all translation-stable, so Proposition~\ref{prop:interp} applies to each: the computational
problems there are self-reducible in Type 3 whatever the covariance obstruction, including
the separating target of Proposition~\ref{prop:sep}, whose degree $3$ needs four queries. The
reduction never rebuilds the covariance defect --- the solver returns it implicitly inside its
own answers --- so no uncomputable monomial ever has to be formed.

What it costs is the regime. The queries must be correlated for the interpolation identity to
hold, so the transfer runs through a union bound and the bound is vacuous unless
$\Succ_{\mathrm{avg}}(\mathcal A)>d/(d+1)$. In the classical setting that threshold is soft:
Berlekamp--Welch decoding brings the tolerable disagreement to any constant below
$1/2$~\cite{GS92}, and list decoding brings it past $1/2$~\cite{Sudan97,STV01}. Neither is
available here. Berlekamp--Welch solves $N(t_i)=y_i\,E(t_i)$ for the coefficients of $N$ and
$E$, a linear system whose \emph{coefficient matrix contains the received values} $y_i$; the
reduction holds only $[y_i]_\star$, and forming that system would require multiplying and
inverting encoded values. The same objection applies to the list-decoding refinements. So the
very criterion that lets interpolation survive the encoding --- linearity in the received word
--- excludes the machinery that would soften its threshold, and $d/(d+1)$ is a genuine wall
rather than an artifact.

\begin{remark}[challenge transport]
\label{rem:decblock}
Interpolation does not reach the decisional problem. Recovering $f(\vx_0)$ from the answers is
linear and hence free in the exponent, but a decisional reduction must also run in the other
direction, transporting its own challenge $Z$ forward to a challenge for each rebuilt
instance, and decision bits do not interpolate. Concretely, for the target
$f=X_1^{2}+X_1X_2^{2}+X_2^{3}$ of Proposition~\ref{prop:sep} the reduction would need the
coefficients of $g(t)=f(\vx_0+t\vb)$; the coefficient of $t$ is the directional derivative
\[
  b_1\bigl(2X_1+X_2^{2}\bigr)+b_2\bigl(2X_1X_2+3X_2^{2}\bigr)
  =2b_1X_1+2b_2X_1X_2+(b_1+3b_2)X_2^{2},
\]
whose $X_2^{2}$ term lies outside $\LL_T$ in Type 3 unless $b_1+3b_2=0$ --- that is, unless
$\vb$ points along the one direction whose orbits are the degenerate lines of
Proposition~\ref{prop:sep}.
\end{remark}

\subsection{What each technique buys}

The two techniques are not ordered by reach. On \emph{computational} problems interpolation is
the more permissive of the two: it asks nothing of $f$ but its degree, so wherever the sources
are translation-stable it succeeds regardless of the covariance obstruction, and
Proposition~\ref{prop:interp} applies to every target in this note. What separates them is the
regime. Theorem~\ref{thm:master} transfers advantage exactly, and reaches the decisional
problem; Proposition~\ref{prop:interp} says nothing below solver advantage $d/(d+1)$ and
nothing about decisional problems at all (Remark~\ref{rem:decblock}). The gap is not one of
tightness --- a reduction
requiring $\Succ_{\mathrm{avg}}>3/4$ is not a loose reduction, it does not exist below that
threshold --- and, by the preceding subsection, it cannot be closed by the classical decoding
machinery, because that machinery is not linear in the received word. Since Uber assumptions
are stated at negligible advantage, and very often in decisional form, the covariance route is
in practice the only one available; that is why Theorem~\ref{thm:master}, and not
interpolation, is the tool used throughout this note, and why the separation of
Proposition~\ref{prop:sep} is meaningful despite Proposition~\ref{prop:interp}.

Two further points of contact are worth recording, with their differences. Shift-based
self-correction in the style of Blum, Luby and Rubinfeld~\cite{BLR93} uses the same
$\vx\mapsto\vx+\vb$ as the CDH and DDH reduction of Section~\ref{sec:nonbilinear}, but their
corrector queries the oracle twice per trial and takes a majority over repetitions, tolerating
a constant error rate, where the reduction here makes one query and computes the correction
itself from the instance, transferring advantage exactly. And the worst-case-to-average-case
reductions of lattice cryptography, from Ajtai~\cite{Ajtai96} to Regev~\cite{Regev09}, share
the motivation of this note --- basing an average-case assumption on a worst-case one --- but
are reductions \emph{between distinct problems} rather than self-reductions, and rest on
lattice geometry rather than on polynomial structure; Regev's is moreover a quantum reduction.

\section{Deciding whether the affine technique applies}
\label{sec:decide}

Conditions \textup{(C1)} and \textup{(C2)} are polynomial conditions on the entries of $\tau$,
and \textup{(C3)} is a statement about the size of an orbit. Finding an admissible $\mathcal T$
is therefore a computation rather than a matter of ingenuity, and the hand arguments of
Sections~\ref{sec:nonbilinear}--\ref{sec:type3} are all instances of one procedure. This
section writes it out and, at each step, names the standard operation of computational algebra
that performs it. Nothing bespoke is needed: the whole of it is subspace arithmetic, one
stabilizer computation, one Jacobian kernel, one ideal of maximal minors, and one primary
decomposition, in that order. What the procedure decides is the applicability of
Theorem~\ref{thm:master}, not random self-reducibility as such;
Section~\ref{sec:limits} is precise about the gap.

\subsection{Set-up}

Fix a setting and an instance $(R,S,T,f;\ch)$ in $n$ variables over $\Zp$, and let $D$ be the
largest degree occurring among $f$ and the generators listed for $\LL_1,\LL_2,\LL_T$ in
Lemma~\ref{lem:comp} --- so $D\ge\deg r+\deg s$ for the pairing products, not merely the
largest degree in $R\cup S\cup T$. Affine substitution does not raise
degree, so every polynomial the procedure meets lies in $P_D\coloneqq\Zp[\vX]_{\le D}$, a
$\Zp$-space of dimension $N=\binom{n+D}{n}$; represent its elements by coefficient vectors.
Treat the $n^{2}+n$ entries of $(A,\vb)$ in $\tau_{A,\vb}\colon\vx\mapsto A\vx+\vb$ as
indeterminates; for $q\in P_D$ the coefficients of $q\circ\tau_{A,\vb}$ are then polynomials of
degree at most $\deg q$ in those indeterminates, obtained by expansion. Composition with an
affine map is $\Zp$-linear on $P_D$ and respects composition, so it is a representation
\[
  \rho\colon\AGL_n(\Zp)\longrightarrow\GL(P_D),\qquad \rho(\tau)q=q\circ\tau,
\]
whose matrix entries are polynomials of degree at most $D$ in $(A,\vb)$. Everything below is an
operation on $\rho$ and on the subspaces of $P_D$ that Lemma~\ref{lem:comp} supplies.

Two structural facts fix the shape of the procedure. First, \textup{(C1)} and \textup{(C2)} cut
out an algebraic subset $\mathcal T\subseteq\AGL_n(\Zp)$ which is a subgroup, and
Lemma~\ref{lem:stab} below identifies it as a stabilizer, which is where the group structure
asserted in Remark~\ref{rem:dichotomy} comes from. Second, with $\mathcal D$ uniform on
$\mathcal T$ the image $\tau(\vx_0)$ is uniform on the orbit $\mathcal T\vx_0$, because
$\tau\mapsto\tau(\vx_0)$ is $|\Stab_{\mathcal T}(\vx_0)|$-to-one onto that orbit. Deciding
\textup{(C3)} is therefore deciding how large the orbits are, and by the Lang--Weil estimate an
orbit of dimension $e$ carries $\Theta(p^{e})$ points, with implied constants depending only on
$n$ and $D$. For $p$ large relative to those, comparing $\dim\mathcal T\vx_0$ with $n$ decides
\textup{(C3)} up to exactly the $O(1/p)$ slack that Remark~\ref{rem:c3} already tolerates.

\subsection{The admissible set is a stabilizer}

Conditions \textup{(C1)} and \textup{(C2)} look like conditions of different kinds --- one a
membership requirement on the rebuilt sources, the other a covariance requirement on the
challenge. Under $\rho$ they are two instances of one condition: each says that $\tau$ fixes
something.

\begin{lemma}[stabilizer form]
\label{lem:stab}
Write $\bar f\coloneqq f+\LL_\star\in P_D/\LL_\star$. For $\tau\in\AGL_n(\Zp)$:
\begin{enumerate}[label=\textup{(\roman*)},leftmargin=2.6em,itemsep=2pt]
  \item \textup{(C1)} holds if and only if $\rho(\tau)$ stabilizes each of $\LL_1$, $\LL_2$ and
        $\LL_T$;
  \item assuming \textup{(C1)}, so that $\rho(\tau)$ descends to a map $\bar\rho(\tau)$ of
        $P_D/\LL_\star$, condition \textup{(C2)} holds if and only if $\bar\rho(\tau)$ fixes the
        line $\Span{\bar f}$, and $\lambda_\tau$ is then the scalar by which it acts there.
\end{enumerate}
Hence the admissible set is a stabilizer,
\[
  \mathcal T=\bigl\{\tau:\rho(\tau)\LL_\gamma=\LL_\gamma\ (\gamma=1,2,T)
  \ \text{ and }\ \bar\rho(\tau)\Span{\bar f}=\Span{\bar f}\bigr\},
\]
the $\rho$-preimage of an intersection of three subspace stabilizers --- maximal parabolic
subgroups of $\GL(P_D)$, or all of $\GL(P_D)$ in the degenerate case $\LL_\gamma=P_D$ ---
with the stabilizer of a point of $\PP(P_D/\LL_\star)$. In particular $\mathcal T$ is a subgroup
of $\AGL_n(\Zp)$, and $\lambda\colon\mathcal T\to\Zpx$, $\tau\mapsto\lambda_\tau$, is a
character of it.
\end{lemma}
\begin{proof}
(i) In each of the four settings the computable spaces of Lemma~\ref{lem:comp} are generated by
$R,S,T$ under multiplication: $\LL_2=\Span{S}$, $\LL_1=\Span{R}$ or $\Span{R\cup S}$, and
$\LL_T=\LL_1\cdot\LL_2+\Span{T}$, with $\LL_1=\LL_2=\Span{R}$ in the single-group settings and
$\LL_T$ absent when there is no pairing. If $\rho(\tau)$ stabilizes the three spaces then
\textup{(C1)} follows, since $R\subseteq\LL_1$, $S\subseteq\LL_2$ and $T\subseteq\LL_T$.
Conversely assume \textup{(C1)}. Then $\rho(\tau)\Span{S}\subseteq\LL_2=\Span{S}$; and both
$\Span{R}$ and $\Span{S}$ land in $\LL_1$ --- in Type 2 because $\LL_2\subseteq\LL_1$, in the
other settings because $\LL_1=\Span{R}$ and only $R$ is involved --- so
$\rho(\tau)\LL_1\subseteq\LL_1$. Composition with $\tau$ is a ring endomorphism of $\Zp[\vX]$,
so $(uv)\circ\tau=(u\circ\tau)(v\circ\tau)$ and therefore
$\rho(\tau)(\LL_1\cdot\LL_2)\subseteq\LL_1\cdot\LL_2$; with $\rho(\tau)\Span{T}\subseteq\LL_T$
from \textup{(C1)} this gives $\rho(\tau)\LL_T\subseteq\LL_T$. Each inclusion is between spaces
of equal finite dimension under an invertible map, hence an equality.

(ii) The identity $f\circ\tau=\lambda_\tau f+h_\tau$ with $h_\tau\in\LL_\star$ says exactly that
the image of $f\circ\tau$ in $P_D/\LL_\star$ is $\lambda_\tau\bar f$, and the requirement
$\lambda_\tau\in\Zpx$ says the image is a \emph{nonzero} multiple, that is, that the line is
fixed rather than merely mapped into itself. The final claims follow because the preimage of an
intersection of subgroups under a homomorphism is a subgroup, and because the scalar by which a
group acts on a fixed line is multiplicative in the group.
\end{proof}

Lemma~\ref{lem:stab} supersedes the direct check in Remark~\ref{rem:dichotomy} that
$\mathcal T$ is closed under composition and inversion: it is a preimage of a subgroup. It also
places $\lambda$ correctly as a character rather than an incidental scalar, which is what makes
$\lambda_\tau\in\Zpx$ compose in the multi-query condition \textup{(C2$'$)}.

\subsection{Orbits: invariants and a rank locus}

By Lemma~\ref{lem:stab}, $\mathcal T$ is an algebraic subgroup of $\AGL_n$ acting on affine
$n$-space, and \textup{(C3)} asks that its orbits be everything. Two standard descriptions of
orbits apply, one generic and one pointwise, and between them they replace the orbit
bookkeeping of Sections~\ref{sec:nonbilinear}--\ref{sec:type3}.

Generically, orbits are the fibres of the invariant map. By Rosenlicht's theorem on generic
quotients there is a dense open subset of $\overline{\Zp}^{\,n}$ on which the orbits of
$\mathcal T$ are exactly the fibres of the map assembled from a finite generating set of the
invariant field, and $\operatorname{trdeg}\overline{\Zp}(\vX)^{\mathcal T}=
n-\max_{\vx}\dim\mathcal T\vx$. So the generic orbit is $n$-dimensional precisely when
$\mathcal T$ admits no non-constant rational invariant. This is the exact shape of the Type-3
obstruction in Proposition~\ref{prop:sep}: there the invariant is $x_1+3x_2$, the orbits are its
level sets, and the ``no escape on an orbit'' step of that proof is a membership test on the
generic fibre.

Pointwise, the failure locus is where the infinitesimal action drops rank. Write
$\agl_n=\gl_n\ltimes\Zp^{n}$ and let $(M,\vvec)$ act on polynomials as the derivation
\[
  \delta_{M,\vvec}=\sum_{i=1}^{n}(M\vX+\vvec)_i\,\frac{\partial}{\partial X_i},
\]
the derivative at $t=0$ of $q\mapsto q\circ(\mathrm{id}+t\,(M,\vvec))$. Differentiating the
conditions of Lemma~\ref{lem:stab} at the identity turns them into a \emph{linear} system in
$(M,\vvec)$:
\begin{enumerate}[label=\textup{(L\arabic*)},leftmargin=3.4em,itemsep=2pt]
  \item $\delta r\in\LL_1$, $\delta s\in\LL_2$ and $\delta t\in\LL_T$ for all $r\in R$,
        $s\in S$, $t\in T$;
  \item $\delta f\in\Span{f}+\LL_\star$, the coefficient of $f$ being $d\lambda(\delta)$.
\end{enumerate}
Its solution space $\widehat\tge$ is the tangent space at the identity to the scheme cut out by
\textup{(C1)} and \textup{(C2)}, so $\Lie\mathcal T\subseteq\widehat\tge$ and
$\dim\mathcal T\le\dim\widehat\tge$, with equality when that scheme is smooth at the identity.

\begin{lemma}[rank locus]
\label{lem:rank}
Let $\tge\subseteq\Lie\mathcal T$ be spanned by $(M_1,\vvec_1),\dots,(M_d,\vvec_d)$ and let
$\Theta(\vX)$ be the $n\times d$ matrix of affine-linear forms whose $i$-th column is
$M_i\vX+\vvec_i$. Then $\rk\Theta(\vx)\le\dim\mathcal T\vx$ for every $\vx$, with equality when
$\tge=\Lie\mathcal T$ and the orbit map $\tau\mapsto\tau(\vx)$ is separable. Consequently the
failure locus $Z\coloneqq\{\vx:\dim\mathcal T\vx<n\}$ satisfies
\[
  Z\ \subseteq\ V\bigl(I_n(\Theta)\bigr),
\]
where $I_n(\Theta)$ is the ideal of $n\times n$ minors of $\Theta$ --- the $0$-th Fitting ideal
of $\operatorname{coker}\Theta$ --- generated in degree at most $n$, and equal to $(0)$ when
$d<n$.
\end{lemma}
\begin{proof}
Orbits of an algebraic group are homogeneous, hence smooth, so
$\dim\mathcal T\vx=\dim T_{\vx}(\mathcal T\vx)$. The orbit map $\mathcal T\to\mathcal T\vx$ has
differential $\mu\mapsto\mu(\vx)$ at the identity, whose image lies in $T_{\vx}(\mathcal T\vx)$
and contains the column span of $\Theta(\vx)$, with equality when $\tge=\Lie\mathcal T$; it is
the whole tangent space exactly when the orbit map is separable. The displayed inclusion is the
rank bound read contrapositively: $\rk\Theta(\vx)=n$ forces $\dim\mathcal T\vx=n$.
\end{proof}

Lemma~\ref{lem:rank} is what makes the failure locus computable rather than merely definable.
The entries of $\Theta$ are affine-linear, so its maximal minors have degree at most $n$ and are
written down by expansion; the locus they cut out is a determinantal variety, and its components
come from a primary decomposition. In every example of this note the inclusion is an equality
and the minors reproduce the loci found by hand: $(1)$ --- empty locus --- for CDH, DDH and
BDDH, the ideal $(x_1)$ for co-CDH, the ideal $(x_1x_2)$ for the monomial targets of
Proposition~\ref{prop:nomono}, and $(0)$ --- the whole space --- for the Type-3 half of
Proposition~\ref{prop:sep}, where $\tge$ is one-dimensional and no $2\times 2$ minor exists.

\subsection{The procedure}

\begin{procedurebox}
\textbf{\textsc{Admissible}}$(\Gamma,R,S,T,f,\ch,p)$ --- decides applicability of
Theorem~\ref{thm:master}. The bracketed name on each step is the operation it calls.\\[2pt]
\emph{Returns} one of \textsc{Not-Hard}, \textsc{Exact}, \textsc{Relaxed},
\textsc{Affine-Fails}, together with $\mathcal T$ and, in the last case, a witness locus.
\begin{enumerate}[label=\textup{\arabic*.},leftmargin=2.2em,itemsep=3pt]
  \item \emph{Computable spaces} \textup{[image of a multiplication map]}. From
        Lemma~\ref{lem:comp} form the generating sets for the setting $\Gamma$ and compute bases
        of $\LL_1,\LL_2,\LL_T\subseteq P_D$ by Gaussian elimination; the target space is the
        image of the multiplication map $\LL_1\otimes\LL_2\to P_D$ plus $\Span{T}$. Put
        $\LL_\star\coloneqq\LL_{\ch}$.
  \item \emph{Hardness} \textup{[subspace membership]}. Test $f\in\LL_\star$. If so, return
        \textsc{Not-Hard}: the challenge is already computable from the instance and there is
        no assumption to reduce.
  \item \emph{Admissible group} \textup{[stabilizer; Gr\"obner basis of a saturated ideal]}. Let
        $\pi\colon P_D\to P_D/\LL_\star$ and $\bar f\coloneqq\pi(f)$, nonzero by step 2. With
        the $n^{2}+n$ entries of $(A,\vb)$ and one further indeterminate $\lambda$, let $I$ be
        generated by
        \begin{enumerate}[label=\textup{(\alph*)},leftmargin=2.2em,topsep=2pt,itemsep=1pt]
          \item the images of $\mathrm{coef}(r\circ\tau)$ in $P_D/\LL_1$, of
                $\mathrm{coef}(s\circ\tau)$ in $P_D/\LL_2$, and of $\mathrm{coef}(t\circ\tau)$
                in $P_D/\LL_T$, for each $r\in R$, $s\in S$, $t\in T$ --- this is
                \textup{(C1)} in the form of Lemma~\ref{lem:stab}(i);
          \item the coefficients of $\pi(f\circ\tau)-\lambda\bar f$ --- this is \textup{(C2)}
                in the form of Lemma~\ref{lem:stab}(ii).
        \end{enumerate}
        Compute a Gr\"obner basis of the saturation $I:(\lambda\det A)^{\infty}$, by
        Rabinowitsch in one extra variable. Its variety is $\mathcal T$, the character
        $\lambda$ is the $\lambda$-coordinate, and $\dim\mathcal T$ is the Hilbert dimension of
        the basis. If $\mathcal T$ is trivial, return \textsc{Affine-Fails} with witness
        $\Zp^{n}$.
  \item \emph{Infinitesimal action} \textup{[Jacobian kernel: one linear system]}. Solve
        \textup{(L1)--(L2)} for $\widehat\tge$. If $\dim\widehat\tge<n$ then
        $\dim\mathcal T\le\dim\widehat\tge<n$, every orbit is deficient, and by
        Remark~\ref{rem:dichotomy} \textup{(C3)} fails everywhere: return
        \textsc{Affine-Fails} with witness $\Zp^{n}$. Otherwise retain those basis elements
        whose one-parameter subgroups can be verified against \textup{(C1)--(C2)} by
        substitution, obtaining $\tge\subseteq\Lie\mathcal T$; retaining fewer only makes the
        verdict more conservative.
  \item \emph{Transitivity} \textup{[the translation subgroup, a $\Zp$-subspace]}. Let
        $W\coloneqq\{\vb:\tau_{\mathrm{id},\vb}\in\mathcal T\}$, a subgroup of $(\Zp^{n},+)$ and
        hence a $\Zp$-subspace; compute it as $\{\vvec:(0,\vvec)\in\widehat\tge\}$ and verify each
        basis vector against \textup{(C1)--(C2)}. If $W=\Zp^{n}$ then $\mathcal T$ is
        transitive: return \textsc{Exact}, and Theorem~\ref{thm:master} applies with unrelaxed
        \textup{(C3)} and exact advantage transfer.
  \item \emph{Failure locus} \textup{[maximal minors: a Fitting ideal, then its minimal
        primes]}. Otherwise form $\Theta$ from $\tge$ and compute $I_n(\Theta)$. By
        Lemma~\ref{lem:rank} the failure locus satisfies $Z\subseteq V(I_n(\Theta))$; take the
        minimal primes $Z_1,\dots,Z_m$ of $V(I_n(\Theta))$ by primary decomposition. If
        $I_n(\Theta)=(1)$ the orbits all have dimension $n$ and, by Remark~\ref{rem:dichotomy},
        return \textsc{Relaxed}. If $I_n(\Theta)=(0)$, return \textsc{Affine-Fails} with witness
        $\Zp^{n}$.
  \item \emph{Escape clause} \textup{[ideal membership; normal form with the field equations
        adjoined]}. For each component $Z_j$ test both:
        \begin{enumerate}[label=\textup{(\alph*)},leftmargin=2.2em,topsep=2pt,itemsep=1pt]
          \item \emph{detectability} --- is the ideal $I(Z_j)$ generated by polynomials lying in
                $\LL_\gamma$ for some group $\gamma$? If so the reduction decides membership by
                testing each $[q(\vx_0)]_\gamma$ against the identity;
          \item \emph{computability} --- is $f\in\LL_\star$ \emph{as a function} on $Z_j(\Zp)$?
                Adjoin the field equations and reduce: compute normal forms modulo a Gr\"obner
                basis of $I(Z_j)+(X_1^{p}-X_1,\dots,X_n^{p}-X_n)$, an ideal that is radical and
                whose variety is exactly $Z_j(\Zp)$, and test membership of the normal form of
                $f$ in the span of those of $\LL_\star$.
        \end{enumerate}
        If both hold for every $j$, return \textsc{Relaxed}: Theorem~\ref{thm:master} applies
        via Remark~\ref{rem:c3}, with advantage transfer up to $O(1/p)$. Otherwise return
        \textsc{Affine-Fails} with the first offending $Z_j$ as witness.
\end{enumerate}
\end{procedurebox}

Three of these steps are worth a word. Step 3 introduces $\lambda$ as an indeterminate rather
than imposing the proportionality $\pi(f\circ\tau)\parallel\bar f$ through the vanishing of the
$2\times 2$ minors of $[\,\pi(f\circ\tau)\mid\bar f\,]$: the two cut out the same set, but the
minors number $\binom{\dim P_D/\LL_\star}{2}$ and have degree $2D$, whereas the linear-in-%
$\lambda$ form has $\dim P_D/\LL_\star$ generators of degree $D$, and the saturation at
$\lambda$ is exactly the requirement $\lambda_\tau\in\Zpx$. Step 5 tests a sufficient condition
for transitivity rather than transitivity itself, because that is all the \textsc{Exact} verdict
needs and because the translation subgroup, being a subgroup of $(\Zp^{n},+)$, is a subspace and
so costs one linear solve; a transitive $\mathcal T$ with $W\neq\Zp^{n}$ is merely reported as
\textsc{Relaxed}, at a cost of $O(1/p)$ in the advantage. Step 7(b) is where the distinction
between polynomials and functions bites --- $f$ may lie outside $\LL_\star$ as a polynomial and
inside it as a function on $Z_j(\Zp)$ --- and the standard device for it is to adjoin the field
equations $X_i^{p}-X_i$, which is what forces the hypothesis $p\ge 5$ in
Proposition~\ref{prop:sep}: at $p=3$ the relation $s^{3}=s$ is among them.

\begin{proposition}[correctness]
\label{prop:decide}
\textsc{Admissible} terminates, and its output is correct: it returns \textsc{Not-Hard} exactly
when $f\in\LL_\star$; it returns \textsc{Exact} only when some efficiently sampleable
$\mathcal D$ satisfies \textup{(C1)--(C3)} as stated in Theorem~\ref{thm:master}; it returns
\textsc{Relaxed} only when Theorem~\ref{thm:master} applies in the form of
Remark~\ref{rem:c3}; when it returns \textsc{Affine-Fails} with witness $\Zp^{n}$, no
distribution on any subgroup of $\AGL_n(\Zp)$ satisfies \textup{(C1)--(C3)} at any point, even
relaxed; and when it returns \textsc{Affine-Fails} with witness a component $Z_j$ on which
$\dim\mathcal T\vx<n$, the same holds at the points of $Z_j$.
\end{proposition}
\begin{proof}
Termination is that every step is a finite linear-algebra or Gr\"obner computation over $\Zp$.
Step 2 is the generic-model hardness criterion of~\cite{BBG05,Boyen08}. Steps 3 and 4 are exact
transcriptions of Lemma~\ref{lem:stab} and of its derivative at the identity: membership of
$q\circ\tau$ in a subspace is the vanishing of its image in the quotient, and
$f\circ\tau=\lambda_\tau f+h_\tau$ with $h_\tau\in\LL_\star$ says that $\pi(f\circ\tau)$ equals
$\lambda\bar f$, the saturation enforcing $\lambda_\tau\neq 0$ and $\det A\neq 0$. Since
$\mathcal T$ is a group and $\mathcal D$ is uniform on it, $\tau(\vx_0)$ is uniform on
$\mathcal T\vx_0$; so \textup{(C3)} holds exactly when the orbit is all of $\Zp^{n}$, which
step 5 certifies through a transitive subgroup of translations, and holds to within $O(1/p)$
when the orbit has full dimension, which step 6 certifies through Lemma~\ref{lem:rank}.
Maximality is what gives the negative half: by Lemma~\ref{lem:stab}, $\mathcal T$ contains
\emph{every} $\tau$ meeting \textup{(C1)} and \textup{(C2)}, so any admissible subgroup is
contained in $\mathcal T$ and its orbits in $\mathcal T$-orbits; if the $\mathcal T$-orbit of a
point of $Z_j$ is already deficient and $f$ is not computable on it, no smaller group can do
better. For the witness $\Zp^{n}$ at step 4 this is unconditional, since
$\dim\mathcal T\le\dim\widehat\tge$ bounds the orbit dimension at every point at once.
\end{proof}

\begin{remark}[the negative verdict at step 7 is conservative]
\label{rem:conservative}
Lemma~\ref{lem:rank} gives $Z\subseteq V(I_n(\Theta))$ and not equality, so a component $Z_j$
returned by step 6 is a \emph{candidate} failure locus: it is a genuine one when
$\tge=\Lie\mathcal T$ and the orbit maps along $Z_j$ are separable, and otherwise the orbits
along $Z_j$ may be larger than $\Theta$ can see, in which case the procedure has failed to
certify rather than certified a failure. The positive verdicts are unaffected, since they rest
on the inclusion in the direction Lemma~\ref{lem:rank} does supply. When the distinction
matters it is settled by the more expensive test the rank locus replaces: compute the orbit of a
generic point of $Z_j$ as the image of the morphism $\tau\mapsto\tau(\vx_0)$, by eliminating
$(A,\vb)$ from $\vy=A\vx_0+\vb$ over $I$ and applying the closure theorem, and compare its
dimension with $n$. In every example of this note the two agree.
\end{remark}

\begin{remark}[cost]
Only step 3 is expensive: its Gr\"obner computation is doubly exponential in $n^{2}+n+1$ in the
worst case, though the systems arising here are far from worst case --- they are the
coefficient-matching equations of Section~\ref{sec:type3}, near-triangular in the entries of
$A$, and for $n=2$ they are seven variables and are solved by hand. It is also, for the verdict,
optional: every one of the four outputs is reached from steps 1, 2 and 4--7, which need only
$\widehat\tge$ and $\tge$. What step 3 adds is $\mathcal T$ in closed form and the character
$\lambda$, and with them the maximality on which the negative half of
Proposition~\ref{prop:decide} rests. Steps 1, 2, 5 and 7 are $O(N^{3})$ linear algebra with
$N=\binom{n+D}{n}$; step 4 is a linear system in $n^{2}+n$ unknowns; step 6 is
$\binom{d}{n}$ determinants of size $n$ followed by a primary decomposition in $n$ variables of
an ideal generated in degree at most $n$. Every primitive named above is a standard call in a
general-purpose system such as Singular, Macaulay2 or Magma.
\end{remark}

\subsection{The primitives, and the procedure on the examples of this note}

Collecting the bracketed names gives the whole of \textsc{Admissible} as a pipeline:

\begin{center}
\small
\begin{tabular}{@{}rll@{}}
\toprule
step & task & operation \\
\midrule
1 & the computable spaces $\LL_\gamma$ & image of a multiplication map \\
2 & hardness, $f\in\LL_\star$ & subspace membership \\
3 & the admissible group $\mathcal T$ & stabilizer of subspaces and of a \\
  & & point of $\PP(P_D/\LL_\star)$; Gr\"obner \\
  & & basis of a saturated ideal \\
3 & the scalars $\lambda_\tau$ & the character of $\mathcal T$ on $\Span{\bar f}$ \\
4 & $\Lie\mathcal T$ & kernel of the Jacobian at the identity \\
5 & transitivity & the translation subspace $W\le\Zp^{n}$ \\
6 & the failure locus & $0$-th Fitting ideal $I_n(\Theta)$, then \\
  & & its minimal primes \\
7 & detectability of $Z_j$ & ideal membership \\
7 & computability on $Z_j$ & normal form modulo
                             $I(Z_j)+(X_i^{p}-X_i)$ \\
\bottomrule
\end{tabular}
\end{center}

\noindent
Run on the assumptions treated above, it reproduces every one of the hand arguments, with the
whole of the orbit analysis coming from a single determinant:

\begin{center}
\footnotesize
\setlength{\tabcolsep}{4.5pt}
\begin{tabular}{@{}llll@{}}
\toprule
target & $\tge$, as vector fields & $I_n(\Theta)$ & verdict \\
\midrule
CDH, DDH (Prop.~\ref{prop:cdh})
  & $\Span{\partial_1,\partial_2,X_1\partial_1,X_2\partial_2}$
  & $(1)$ & \textsc{Exact} \\
BDDH (Prop.~\ref{prop:bddh})
  & $\Span{\partial_i,\ X_i\partial_i:i\le 3}$
  & $(1)$ & \textsc{Exact} \\
co-CDH (Prop.~\ref{prop:cocdh})
  & $\Span{X_1\partial_1,X_2\partial_2,\partial_2}$
  & $(x_1)$ & \textsc{Relaxed} \\
$X_1^{i}X_2^{j}$ (Prop.~\ref{prop:nomono})
  & $\supseteq\Span{X_1\partial_1,X_2\partial_2}$
  & $\ni x_1x_2$ & \textsc{Relaxed} \\
Prop.~\ref{prop:sep}, Type 2
  & $\Span{\partial_1,\partial_2,(2X_1{+}X_2)\partial_1+X_2\partial_2}$
  & $(1)$ & \textsc{Exact} \\
Prop.~\ref{prop:sep}, Type 3
  & $\Span{-3\partial_1+\partial_2}$
  & $(0)$ & \textsc{Affine-Fails} \\
\bottomrule
\end{tabular}
\end{center}

\noindent
The last two rows are the separation, and in this form it is a statement about one number: the
Type-2 conditions leave a three-dimensional $\tge$ containing both translations, so
$\Theta$ has an invertible $2\times 2$ submatrix and the failure locus is empty, while the
Type-3 conditions cut $\tge$ down to a single vector field, $\Theta$ becomes a $2\times 1$
matrix, and no $2\times 2$ minor exists to be nonzero. The intermediate rows show the two ways a
non-empty failure locus can still be harmless, both settled by step 7: the locus $x_1=0$ for
co-CDH and $x_1x_2=0$ for the monomial targets are detectable by an equality test on the
encodings, and $f$ vanishes on each, so the reduction answers there without an oracle call.
Row three also shows what maximality buys: the group the procedure returns for co-CDH is
three-dimensional, strictly larger than the $\{(\alpha x_1,\,x_2+b_2)\}$ used in the proof of
Proposition~\ref{prop:cocdh}, which chose a sufficient subgroup rather than the largest one.

\subsection{What the procedure does not decide}
\label{sec:limits}

\textsc{Affine-Fails} is not a proof that the assumption fails to be random self-reducible. It
is a proof about one technique, and three gaps remain open beneath it.

First, multi-query reductions. Proposition~\ref{prop:interp} solves the computational problem
whenever the sources are translation-stable and $\deg f\le p-2$, whatever the verdict of
\textsc{Admissible} --- as it does for the separating target of Proposition~\ref{prop:sep}. A
run returning \textsc{Affine-Fails} should therefore be followed by that test, which is a
one-line check. The general linear multi-query condition \textup{(C2$'$)} is decidable too, but
by a different route: it asks whether $f$ lies in
$\Span{f\circ\tau:\tau\ \text{satisfies (C1)}}+\LL_\star$. By Lemma~\ref{lem:stab}(i) the
$\tau$ satisfying \textup{(C1)} form a group $\mathcal G_1$ --- the same stabilizer as
$\mathcal T$, without the condition on $\bar f$ --- so that span is the $\rho(\mathcal
G_1)$-submodule of $P_D$ generated by $f$, and step 3 run without its part (b) returns
$\mathcal G_1$. The test is then the standard closure computation: apply generators of
$\mathcal G_1$ to a spanning set repeatedly until the dimension stabilises, which happens after
at most $N$ rounds, and test membership of $f$ in the result plus $\LL_\star$. It is necessary
but not sufficient, since \textup{(C2$'$)} also requires the $\tau_i$ realizing the combination
to be marginally distributed as in \textup{(C3)}.

Second, non-linear recombination. Feigenbaum and Fortnow~\cite{FF93} allow the answers to be
combined by an arbitrary polynomial-time function, and nothing here decides that condition; the
procedure sees only the $\Zp$-linear fragment that the group encoding makes free.

Third, reparametrizations outside $\AGL_n(\Zp)$, and reductions that are not reparametrization
based at all. A genuine impossibility result for the decisional problem would need a lower-bound
argument such as a meta-reduction, as noted after Proposition~\ref{prop:sep}.

\textsc{Admissible} is accordingly a decision procedure for a sufficient condition, and a
semi-decision procedure for random self-reducibility: it can certify it, and it cannot refute
it.

\section{Discussion}
\label{sec:discussion}

Across the four settings the reduction of Theorem~\ref{thm:master} is one and the same;
only the admissibility conditions move, and they move only because the computable target
space $\LL_\star$ grows with the available products:
\[
  \text{none}\ \subsetneq\ R\cdot S\ (\text{Type 3})\ \subsetneq\ R\cdot S\cup S\cdot S\
  (\text{Type 2})\ \subseteq\ R\cdot R\ (\text{Type 1, single group}).
\]
A larger $\LL_\star$ makes the covariance condition (C2) easier to meet, so the
affine-self-reducible subclass is largest in Type 1 and smallest in Type 3, with the
non-bilinear case the degenerate product-free floor. Proposition~\ref{prop:sep} witnesses the
strict gap between Type 2 and Type 3, at a target where the Type-3 reduction fails at every
worst-case point --- the most the failure locus can ever be, by
Remark~\ref{rem:dichotomy}. The gap is one for decisional reductions and for reductions that
must work at negligible advantage: Proposition~\ref{prop:interp} still solves the
computational problem in Type 3 at high solver advantage, as it does for every target here.
Proposition~\ref{prop:nomono} marks the other boundary, that no monomial target can separate
any two settings at all.

Self-reducibility is orthogonal to hardness. The master theorem
of~\cite{BBG05,Boyen08} makes an assumption hard exactly when $f\notin\LL_\star$, whereas
Theorem~\ref{thm:master} needs the defect $f\circ\tau-\lambda_\tau f$ to lie \emph{in}
$\LL_\star$. These constrain different objects, so hard yet self-reducible assumptions
abound: CDH and DDH (non-bilinear), BDDH (Type 1), and co-CDH and every monomial target
(Type 3, Proposition~\ref{prop:nomono}) are all of this kind. Quantitatively, Bauer, Fuchsbauer, and
Loss~\cite{BFL20} show that in the algebraic group model every member of the Uber family is
implied by the $q$-discrete-logarithm assumption, with $q$ governed by the degrees of the
defining polynomials, and Rotem~\cite{Rotem22} refines $q$ to the maximal degree in which a
single variable occurs and gives improved bilinear reductions.

\section{Bibliography}

\begin{conjurabibliography}{EHKRV13}

\bibitem[TW87]{TW87}
M.~Tompa and H.~Woll.
Random self-reducibility and zero knowledge interactive proofs of possession of
information.
In \emph{28th Annual Symposium on Foundations of Computer Science (FOCS 1987)},
pages 472--482. IEEE, 1987.

\bibitem[BF90]{BF90}
D.~Beaver and J.~Feigenbaum.
Hiding instances in multioracle queries.
In \emph{7th Annual Symposium on Theoretical Aspects of Computer Science (STACS 1990)},
volume 415 of \emph{Lecture Notes in Computer Science}, pages 37--48. Springer, 1990.

\bibitem[Lipton91]{Lipton91}
R.~J.~Lipton.
New directions in testing.
In \emph{Distributed Computing and Cryptography}, volume 2 of \emph{DIMACS Series in
Discrete Mathematics and Theoretical Computer Science}, pages 191--202. AMS, 1991.

\bibitem[GS92]{GS92}
P.~Gemmell and M.~Sudan.
Highly resilient correctors for polynomials.
\emph{Information Processing Letters}, 43(4):169--174, 1992.

\bibitem[BLR93]{BLR93}
M.~Blum, M.~Luby, and R.~Rubinfeld.
Self-testing/correcting with applications to numerical problems.
\emph{Journal of Computer and System Sciences}, 47(3):549--595, 1993.
Preliminary version in \emph{STOC 1990}, pages 73--83.

\bibitem[FF93]{FF93}
J.~Feigenbaum and L.~Fortnow.
Random-self-reducibility of complete sets.
\emph{SIAM Journal on Computing}, 22(5):994--1005, 1993.

\bibitem[Ajtai96]{Ajtai96}
M.~Ajtai.
Generating hard instances of lattice problems.
In \emph{28th Annual ACM Symposium on Theory of Computing (STOC 1996)}, pages 99--108.
ACM, 1996.

\bibitem[Sudan97]{Sudan97}
M.~Sudan.
Decoding of Reed--Solomon codes beyond the error-correction bound.
\emph{Journal of Complexity}, 13(1):180--193, 1997.

\bibitem[STV01]{STV01}
M.~Sudan, L.~Trevisan, and S.~Vadhan.
Pseudorandom generators without the XOR lemma.
\emph{Journal of Computer and System Sciences}, 62(2):236--266, 2001.

\bibitem[BBG05]{BBG05}
D.~Boneh, X.~Boyen, and E.-J.~Goh.
Hierarchical identity based encryption with constant size ciphertext.
In \emph{Advances in Cryptology -- EUROCRYPT 2005},
volume 3494 of \emph{Lecture Notes in Computer Science}, pages 440--456. Springer, 2005.

\bibitem[BT06]{BT06}
A.~Bogdanov and L.~Trevisan.
Average-case complexity.
\emph{Foundations and Trends in Theoretical Computer Science}, 2(1):1--106, 2006.

\bibitem[Boyen08]{Boyen08}
X.~Boyen.
The uber-assumption family: a unified complexity framework for bilinear groups.
In \emph{Pairing-Based Cryptography -- Pairing 2008},
volume 5209 of \emph{Lecture Notes in Computer Science}, pages 39--56. Springer, 2008.

\bibitem[GPS08]{GPS08}
S.~D.~Galbraith, K.~G.~Paterson, and N.~P.~Smart.
Pairings for cryptographers.
\emph{Discrete Applied Mathematics}, 156(16):3113--3121, 2008.

\bibitem[Regev09]{Regev09}
O.~Regev.
On lattices, learning with errors, random linear codes, and cryptography.
\emph{Journal of the ACM}, 56(6):article 34, 2009.
Preliminary version in \emph{STOC 2005}, pages 84--93.

\bibitem[EHKRV13]{EHKRV13}
A.~Escala, G.~Herold, E.~Kiltz, C.~R\`afols, and J.~Villar.
An algebraic framework for Diffie-Hellman assumptions.
In \emph{Advances in Cryptology -- CRYPTO 2013},
volume 8043 of \emph{Lecture Notes in Computer Science}, pages 129--147. Springer, 2013.
Journal version: \emph{Journal of Cryptology} 30(1):242--288, 2017.

\bibitem[BFL20]{BFL20}
B.~Bauer, G.~Fuchsbauer, and J.~Loss.
A classification of computational assumptions in the algebraic group model.
In \emph{Advances in Cryptology -- CRYPTO 2020},
volume 12171 of \emph{Lecture Notes in Computer Science}, pages 121--151. Springer, 2020.

\bibitem[Rotem22]{Rotem22}
L.~Rotem.
Revisiting the uber assumption in the algebraic group model: fine-grained bounds in
hidden-order groups and improved reductions in bilinear groups.
In \emph{3rd Conference on Information-Theoretic Cryptography (ITC 2022)},
volume 230 of \emph{LIPIcs}, article 13. Schloss Dagstuhl, 2022.

\end{conjurabibliography}

\end{document}
